% -----------------------------------------------------------------------------
% This file is automatically generated.
% Do not edit manually.
% Source: notebooks/support/auxiliary-variables-and-stratification/support.Rmd
% -----------------------------------------------------------------------------

\par\addvspace{\topsep}
\begingroup
\fvset{listparameters={%
  \setlength{\topsep}{0pt}%
  \setlength{\partopsep}{0pt}%
  \setlength{\parsep}{0pt}%
  \setlength{\itemsep}{0pt}%
}}
\begin{Verbatim}[breaklines=true,formatcom=\color{ada_blue}]
arSample <- evaluate(arSample, av = audit_values$av2)
arSample$evalResults$Estimates
\end{Verbatim}
\begin{Verbatim}[breaklines=true,formatcom=\color{ada_light_blue}]
                  mpu       diff      ratio       regr
audit value  10873416 13306202.0 13263601.5 13285520.3
misstatement  2626584   193798.0   236398.5   214479.7
precision     2100023   533347.2   551006.0   532846.0
lower bound   8773393 12772854.8 12712595.5 12752674.2
upper bound  12973439 13839549.2 13814607.5 13818366.3
effective df       99        5.0        5.0        5.0
\end{Verbatim}
\endgroup
\par\addvspace{\topsep}
\noindent The number of differences is stored in the \ttblue{\#\_Errors} attribute, under each of the estimation methods. This way, you can tell if precision was calculated with the sporadic error (k = 4 to 19) method, or using the standard CVS method.
